\section*{\centering\zihao{3}结\quad 论}
\addcontentsline{toc}{section}{结\quad 论}

本文设计并实现了一套部署在树莓派设备上的人脸识别考勤系统。该系统面向真实考勤场景，在有限算力的嵌入式设备上完成了从人脸注册、活体挑战、身份识别到自动签到的端到端业务流程。

本文完成的主要工作如下。

第一，完成了人脸注册流的设计与实现。采用分步动作挑战机制替代传统的单张照片上传，通过正脸、侧转、再正脸三步抓拍结合同人一致性校验，有效提高了注册照的真实性。设计了可独立扩展的质量校验流水线，涵盖单人检测、人脸面积、模糊度、亮度、姿态和翻拍检测六个校验器，每个校验器遵循开闭原则独立实现。引入审核机制，新上传人脸档案需经管理员审核通过后才参与终端识别。

第二，完成了实时识别考勤流的设计与实现。设计了包含单人门禁、人脸面积门禁、编码匹配门禁、多帧稳定确认、考勤时间窗和冷却时间六道约束的严格识别引擎，确保签到结果的可靠性。引入OpenVINO anti-spoof-mn3模型进行被动活体前置判断，采用三区间策略（放行/灰区/低分连续确认）降低真人误杀率。设计了2秒滑动窗口多帧融合机制，对real\_score进行平滑处理，并辅以屏幕重放风险信号分析，实现了不要求用户做配合动作的无感活体检测。

第三，完成了识别后端的优化。采用OpenCV YuNet检测加SFace特征提取的识别链路，通过配置项控制后端选择。实测表明YuNet+SFace方案在树莓派上的单帧处理速度满足实时性要求，平均处理耗时约136毫秒。

第四，完成了系统架构与接口设计。采用五层架构（配置层、数据层、摄像头采集层、Web页面层、识别业务层），层与层之间职责清晰。实现了基于FastAPI的外网用户门户和设备侧管理员API，Windows端通过Tailscale网络远程访问树莓派接口，实现了人脸审核、成员管理和考勤记录查询。

本文的系统在以下几个方面体现了一定的创新性。动作挑战活体检测将注册端的活体验证从``信任一张照片''升级为``信任一个动作过程''，通过随机生成的动作序列和同人一致性校验降低了预录视频攻击的成功率。多帧被动活体融合通过短时间滑动窗口对活体分数进行平滑，结合屏幕重放风险信号，在不要求用户转头、眨眼等配合动作的前提下实现了对照片和屏幕攻击的风险感知。严格识别多门禁设计将识别从``匹配到就签到''简化逻辑升级为多道约束的精细化判断，降低了误识别和重复签到的概率。采集与识别线程解耦使视频流帧率从2--3FPS提升至12--15FPS，终端画面流畅度显著改善。

系统目前仍存在以下不足。活体检测的屏幕重放风险信号基于轻量图像特征，对高质量手机屏幕的区分能力有限，后续可以考虑引入更复杂的时序分析方法。识别后端在大规模人脸库（千人以上）场景下的检索速度有待优化，后续可以考虑引入FAISS等向量索引库。当前系统仅支持单台树莓派设备，后续可以扩展为多设备协同的考勤网络，实现跨地点的统一考勤管理。用户门户页面仍有原型性质，后续可以根据实际使用反馈优化交互体验和视觉设计。Windows管理端的审核功能和操作日志仍需进一步丰富。

未来工作可以沿着以下方向推进。接入更强大的被动活体检测模型，利用多帧时序信息和rPPG生理信号进一步提升活体判断的准确率。引入考勤规则引擎，支持不同场景下的差异化考勤策略配置。增加数据导出和统计报表功能，为管理人员提供更直观的数据分析工具。探索模型量化和TensorRT加速等技术，进一步降低模型推理的计算开销，提升系统在更低配置设备上的运行能力。
